含截距的回归可先中心化，再用协方差或正交投影求系数。新增解释变量时，把响应与新增变量对旧设计同时残差化，可将多元问题化为一维问题。设计满秩与残差方向非零，分别保证所求系数可识别。

测量误差、样本筛选会改变协方差结构；先写清观测模型和条件总体，再区分线性最优预测与条件均值预测。工具变量须与目标信号相关、与相应误差不相关。投影恒等式属于代数结论，而有限样本的精确 $t$ 检验还需要正确模型下独立、同方差的正态误差。
\par\smallskip
